\section{Reconstruction and versioned evidence}
\label{app:reproduction}

Reproducibility is organized around the relationship between retained inputs, deterministic transformations and named outputs. Input fingerprints bind a calculation to the files it used. Source registries identify the interfaces and historical boundaries, while result manifests connect tables and figures to the relevant processing rules. A successful replay establishes consistency for those inputs and transformations.

The tutorial provides one bounded example per branch, checking retained inputs, native units and clocks. Fresh historical queries depend on archive retention and provider behaviour. Retrieving an individual receipt or block does not establish continuous acquisition across the registered interval.

The reproduction record documents commands, environment versions, validation and acquisition limits. Earlier model experiments remain separate from the revised feature definitions. The catalogue and checksums identify this revision's materials and distinguish them from previous public deposits.
